\titleformat{\chapter}[display]
  {\normalfont\bfseries}{}{0pt}{\Large} % Este comando es para eliminar la aparición de "Capítulo N" en el PDF. Para emplear el formato de la tesis simplemente eliminamos la línea. 
\chapter{Justificación}
Los fulerenos, reconocidos por sus propiedades estructurales y electrónicas únicas, han captado un interés creciente en el ámbito de las aplicaciones biomédicas. Sin embargo, su baja solubilidad en entornos acuosos y fisiológicos constituye un obstáculo crítico que limita significativamente su aplicación directa y su potencial terapéutico.

Para superar esta restricción, la funcionalización química de la superficie del fulereno mediante la adición de aminoácidos, como la glicina, representa una estrategia eficaz. La incorporación de grupos polares puede favorecer la interacción con el agua y modificar el reconocimiento de otras moléculas. Estos efectos dependen de la estructura del aducto y requieren evaluarse por separado.

En este contexto, la implementación de herramientas computacionales de modelado, reactividad y acoplamiento molecular proporciona una comprensión detallada a nivel atómico de los mecanismos de interacción del sistema. De este modo, el presente trabajo contribuye al diseño racional de aductos funcionalizados dirigidos específicamente contra la proteasa NS2B-NS3 del virus del dengue, ofreciendo bases teóricas para el desarrollo de nuevos agentes con potencial biomédico.